\chapter{INTRODUCTION}

\section{Background and Context}

The modern technical recruitment landscape relies heavily on automated screening platforms to process large volumes of candidate applications \cite{ref1}. Resumes serve as the primary initial screening document, summarizing a candidate's technical skills, professional experience, project achievements, and academic background \cite{ref2}. However, traditional resume evaluation pipelines face severe structural challenges when assessing technical competency in software engineering, data science, and DevOps domains \cite{ref3, ref4}. Human recruiters, particularly those operating in high-volume talent acquisition pipelines, often lack the domain-specific technical knowledge required to evaluate complex software architectures, framework proficiencies, or algorithmic experience \cite{ref5}. Consequently, recruitment teams rely heavily on automated Applicant Tracking Systems (ATS) to filter candidate pools prior to human review \cite{ref6, ref7}.

Traditional ATS architectures evaluate resumes through deterministic string matching, regular expressions, or basic term frequency heuristics (e.g., TF-IDF) \cite{ref1, ref8}. While effective at reducing applicant volume, keyword-centric screening introduces severe systemic vulnerabilities \cite{ref9}. Candidates who manipulate their resumes by stuffing technical keywords, listing extensive technology stacks, or embedding invisible text blocks can easily bypass initial automated filters \cite{ref10}. Conversely, highly qualified software engineers who describe their project accomplishments using domain-specific synonyms or architectural descriptions without explicitly repeating buzzwords are frequently misclassified and rejected \cite{ref11, ref12}.

This fundamental reliance on keyword matching has created strong financial and professional incentives for candidate skill exaggeration and resume inflation \cite{ref9, ref13}. Recent empirical studies in technical recruitment indicate that up to 70\% of software engineering resumes contain exaggerated or unsubstantiated skill claims \cite{ref10, ref14}. These misrepresentations range from minor exaggerations of framework proficiency—such as declaring expertise in a library after completing a brief online tutorial—to outright fabrication of programming language experience and project contributions \cite{ref15, ref16}. The rapid proliferation of commercial Generative Artificial Intelligence (AI) tools, such as ChatGPT and specialized resume builders, has further amplified this challenge \cite{ref17, ref18}. Applicants can now instantly generate keyword-optimized, tailored resumes designed specifically to score high on ATS algorithms without possessing the underlying implementation capabilities \cite{ref5, ref19}.

Consequently, organizations incur substantial financial, operational, and cultural costs associated with evaluating unqualified candidates \cite{ref20}. Technical interview panels—typically composed of senior software engineers and engineering managers—must spend valuable development hours conducting technical assessments, take-home code reviews, and live coding interviews for candidates whose actual skills fall far short of their resume claims \cite{ref4, ref21}. When unverified hiring decisions occur, companies face project delays, code quality degradation, increased onboarding overhead, and premature employee turnover \cite{ref14, ref22}. Verifying claimed technical skills empirically before committing engineering bandwidth to technical interview panels represents a critical operational requirement for modern hiring pipelines \cite{ref12, ref23}.

\section{The Problem of Resume Fraud in Technical Hiring}

Resume fraud and skill inflation in technical hiring span multiple levels of severity, each posing distinct challenges for talent acquisition teams \cite{ref9, ref10}:

\begin{enumerate}[leftmargin=*, label=\arabic*.]
    \item \textbf{Superficial Technology Listing:} Candidates routinely list frameworks, database engines, and cloud platforms after minimal exposure, such as cloning a starter repository or reading documentation, creating a severe disconnect between claimed proficiency and production capability \cite{ref11, ref12}.
    \item \textbf{Role and Contribution Inflation:} In collaborative open-source or team projects, candidates frequently claim full architectural ownership of complex software modules when their actual contribution was limited to minor bug fixes or documentation edits \cite{ref15, ref16, ref24}.
    \item \textbf{Generative Resume Fabrication:} Leveraging LLMs to synthesize realistic project descriptions populated with advanced industry jargon, making fabricated experience indistinguishable from genuine project work during initial textual screening \cite{ref5, ref17}.
\end{enumerate}

Traditional verification mechanisms attempt to combat these issues through manual code reviews, take-home programming assignments, or live algorithmic coding platforms \cite{ref4, ref15}. However, each of these conventional approaches introduces substantial recruiter friction, candidate drop-off, or evaluation misalignment \cite{ref21}. Take-home projects require significant candidate time investment (often 5 to 15 hours), leading experienced senior engineers to drop out of recruitment funnels. Furthermore, scoring take-home projects requires dedicated engineering bandwidth, limiting scalability \cite{ref16, ref25}.

Conversely, automated competitive coding platforms measure algorithmic problem-solving under artificial time constraints in isolated sandbox environments \cite{ref15}. While useful for assessing basic data structures and algorithms, competitive coding tests fail to evaluate real-world software engineering competencies, such as framework integration, repository organization, database schema design, asynchronous API handling, continuous integration workflows, and long-term codebase maintenance \cite{ref3, ref26}. A candidate may excel at reversing a binary tree on a competitive platform while lacking the practical skill to construct a RESTful API in Node.js or configure a Dockerized deployment pipeline \cite{ref27, ref28}.

\section{Limitations of Existing Automated Approaches}

Recent research in Natural Language Processing (NLP) and Artificial Intelligence has explored automated resume screening and candidate evaluation \cite{ref1, ref11, ref23}. However, current state-of-the-art technologies exhibit four fundamental structural limitations:

\begin{enumerate}[leftmargin=*, label=\arabic*.]
    \item \textbf{Keyword Matching without Verification:} Traditional ATS platforms score resumes based on keyword frequency, rewarding candidates who overload their resumes with technical terms regardless of actual codebase evidence \cite{ref1, ref8}.
    \item \textbf{Monolithic LLM Scoring Risks:} Recent LLM-based resume screeners prompt a single model to score resume quality. These approaches are prone to hallucination, subjective scoring variance, and an inability to inspect external evidence \cite{ref20, ref17, ref29}.
    \item \textbf{Lack of Codebase Grounding:} Existing systems evaluate resumes in isolation as self-contained textual documents, without cross-referencing claims against external artifacts such as public software repositories \cite{ref30, ref31, ref32}.
    \item \textbf{Absence of Multi-Tenant Data Isolation:} Automated scoring platforms frequently lack candidate data segregation mechanisms, creating risks of cross-candidate data contamination during vector retrieval \cite{ref33, ref34}.
\end{enumerate}

Bridging this gap requires an automated system capable of extracting claimed technical skills from resumes and cross-referencing those claims against concrete, functional source code evidence stored in public version control repositories \cite{ref12, ref7, ref35}.

\section{Motivation for the Study}

Public software repositories on platforms like GitHub provide a rich, empirical record of a candidate's actual software engineering contributions \cite{ref15, ref16, ref24}. Public repositories contain source code files, commit histories, configuration files, directory structures, and project architectures that reflect real-world implementation experience \cite{ref6, ref25}. However, manually inspecting candidate repositories is labor-intensive, requiring senior engineers to clone projects, parse directory structures, evaluate code quality across multiple programming languages, and check dependency manifests \cite{ref25, ref26, ref36}.

Advances in Multi-Agent Large Language Model (LLM) orchestration and Retrieval-Augmented Generation (RAG) offer a promising foundation for automating code-level claim verification \cite{ref30, ref37, ref38}. By decomposing the verification workflow into specialized autonomous agents—dedicated to claim extraction, repository ingestion, semantic vector indexing, and evidence-grounded judgment—technical skill claims can be evaluated objectively against actual source code evidence \cite{ref10, ref27, ref33, ref39}.

\section{Problem Statement}

Formally, given a candidate's PDF resume $R$ containing a set of claimed technical skills $S = \{s_1, s_2, \dots, s_n\}$ and a public GitHub profile handle $G$ associated with a set of repositories $V = \{v_1, v_2, \dots, v_m\}$, the objective is to construct an automated multi-agent system that evaluates each skill claim $s_i \in S$ against the ingested codebase evidence $V$, producing a verified verdict $Y_i \in \{\text{Verified}, \text{Unsubstantiated}\}$ accompanied by an evidence-grounded explanatory rationale $E_i$.

Addressing this problem requires solving four key sub-challenges:
\begin{enumerate}[leftmargin=*, label=\arabic*.]
    \item \textbf{Claim Extraction:} Accurately parsing PDF resume structures to extract distinct technical skills while filtering non-technical context \cite{ref9, ref40}.
    \item \textbf{Code Ingestion \& Preprocessing:} Asynchronously ingesting multi-language source code files while filtering non-informative configuration metadata and auto-generated code \cite{ref12, ref35}.
    \item \textbf{Semantic Indexing \& Isolation:} Vectorizing code chunks into a unified semantic space while enforcing strict multi-tenant candidate data isolation \cite{ref31, ref33, ref41}.
    \item \textbf{Evidence-Grounded Judgment:} Retrieving relevant code contexts to perform chain-of-thought reasoning, eliminating LLM hallucination during verification \cite{ref10, ref11, ref25, ref42}.
\end{enumerate}

\section{Objectives of the Project}

The primary goal of this project is to design, implement, and evaluate \textbf{CodeAudit AI}, a multi-agent LLM framework for authenticating technical resume claims via semantic codebase vectorization. The specific objectives are:

\begin{enumerate}[leftmargin=*, label=\arabic*.]
    \item To design a 4-phase modular multi-agent pipeline comprising Claim Extraction, Code Ingestion, RAG Indexing, and LLM Judgment agents \cite{ref10, ref7, ref37, ref38, ref43}.
    \item To implement a PDF claim extraction engine using PyMuPDF and regular expression boundary matching for 17 technical skill categories \cite{ref40}.
    \item To build an asynchronous repository ingestion module utilizing the GitHub REST API v3 and aiohttp for multi-language source code retrieval \cite{ref12, ref35}.
    \item To construct a local vector indexing pipeline using sentence-transformers (all-MiniLM-L6-v2) embeddings and ChromaDB for metadata-isolated vector storage \cite{ref31, ref33, ref41}.
    \item To formulate a Chain-of-Thought (CoT) prompt template for an LLM Judge Agent to evaluate retrieved code chunks and generate deterministic verdicts \cite{ref10, ref11, ref25, ref42}.
    \item To construct a 30-case benchmark dataset across 5 technical categories and conduct ablation studies to quantify system accuracy, precision, recall, F1-score, and specificity \cite{ref20, ref8, ref23}.
\end{enumerate}

\section{Scope of the Project}

The scope of CodeAudit AI focuses on validating claimed technical skills against public GitHub source code repositories. The system supports 17 technical keywords across programming languages (Python, JavaScript, Java, C++), web frameworks (React, Angular, Django, Flask, Node.js), machine learning libraries (TensorFlow, PyTorch, Scikit-learn), database engines (SQL, MongoDB, PostgreSQL), and DevOps tools (Docker, Kubernetes, Git). Source code ingestion is restricted to public GitHub repositories under a 100KB per-file size cap to prevent resource exhaustion from large binary blobs or auto-generated lockfiles.

The project excludes private repositories requiring organization-level credentials, employment duration verification, soft skills evaluation, and direct automated code refactoring. The system is designed as an objective, empirical pre-screening tool to assist human recruiters and technical interviewers.

\section{Organization of the Report}

This report is organized into five comprehensive chapters:
\begin{enumerate}[leftmargin=*, label={Chapter \arabic*:}]
    \item \textbf{Introduction:} Outlines the background, problem domain, motivation, formal problem statement, project objectives, scope, and organization.
    \item \textbf{Literature Review:} Reviews prior work in automated resume screening, RAG architectures, sentence embeddings, code LLMs, multi-agent frameworks, and ethical AI hiring.
    \item \textbf{Proposed Methodology:} Details the 4-phase system architecture, agent workflows, vector indexing, evaluation metrics equations, technology stack, and Streamlit deployment.
    \item \textbf{Results and Discussion:} Presents experimental setup, benchmark dataset, ablation study findings, category-wise performance, error analysis, latency evaluations, and threats to validity.
    \item \textbf{Conclusion and Future Scope:} Summarizes key contributions, system limitations, future research directions, ethical implications, and concluding remarks.
\end{enumerate}